\chapter{Amputation}
\index{Amputation}

\section*{Definition and Overview}
Amputation is the surgical removal of a limb, or part of a limb, usually of the leg or arm. It is performed when a limb cannot be saved after severe injury or disease, when retaining it threatens life through infection or critically impaired blood supply, or when removal offers a better overall quality of life than preservation. Amputations range from the loss of a digit to major procedures such as below-knee, above-knee, below-elbow, or above-elbow amputation. The decision to amputate is major and is never taken lightly; it follows a thorough assessment of whether the limb can be salvaged. Advances in surgical technique, prosthetics, and rehabilitation mean that many people live active, independent lives after amputation.

\section*{Epidemiology}
Most amputations involve the lower limb. Around 1 in 1,000 to 1 in 500 people in high-income countries live with a major limb amputation, and the number is rising as populations age and the prevalence of diabetes increases. Peripheral artery disease and diabetes account for the majority of lower-limb amputations, particularly in people over the age of 60. Trauma is the leading cause in younger people, and in conflict settings blast injuries predominate. Men undergo amputation about twice as often as women, reflecting the higher prevalence of vascular disease and traumatic injury in men. Significant geographical and ethnic variation exists, largely mirroring rates of diabetes and vascular disease.

\section*{Aetiology and Pathophysiology}
Amputation is undertaken when the limb cannot be saved or when keeping it would endanger health or life:

\begin{itemize}
    \item \textbf{Poor blood supply:} peripheral artery disease, driven by smoking, diabetes, high blood pressure, and high cholesterol, progressively starves the limb of oxygen.
    \item \textbf{Diabetes:} neuropathy and impaired circulation lead to ulcers, infection, and tissue death (gangrene).
    \item \textbf{Trauma:} severe crush, blast, or avulsion injuries that cannot be reconstructed.
    \item \textbf{Infection:} overwhelming or deep infections, including gas gangrene and bone infection (osteomyelitis), that do not respond to treatment.
    \item \textbf{Tumours:} bone or soft tissue cancers may require amputation when limb-sparing surgery is not feasible.
    \item \textbf{Congenital conditions:} a limb may be absent, underdeveloped, or malformed from birth.
    \item \textbf{Frostbite and thermal injury:} severe cold or burns can destroy tissue irreversibly.
    \item \textbf{Smoking:} the single most important modifiable risk factor for the vascular causes of amputation.
\end{itemize}

\section*{Clinical Features}
The features that lead to amputation usually develop over months or years, although trauma presents suddenly:

\begin{itemize}
    \item Non-healing wounds or ulcers on the foot or leg.
    \item Persistent pain, coldness, or colour change in the limb, which may be pale, blue, or black.
    \item Gangrene: dead, blackened, often foul-smelling tissue.
    \item Severe infection with fever and spreading redness.
    \item A limb that cannot be moved or used after injury.
    \item After the operation: pain, swelling, and stiffness at the stump, which normally improve with time and rehabilitation.
\end{itemize}

\section*{Differential Diagnosis}
Amputation is a major decision, so the surgical team first excludes reversible causes and tests whether the limb can be saved:

\begin{itemize}
    \item Acute arterial blockage or compartment syndrome that might be reversed with urgent intervention.
    \item Deep or bone infection that might respond to antibiotics and limited bone debridement alone.
    \item Non-ischaemic ulcers, such as those from venous disease or pressure.
    \item Charcot arthropathy in diabetes, which causes a deformed, swollen foot that is not gangrenous.
    \item Whether angioplasty or bypass surgery could restore adequate blood flow.
\end{itemize}

\section*{When to Seek Medical Care}
Seek urgent medical review for any new pain, colour change, or coldness in a limb, and for any diabetic foot wound that does not heal, because early treatment may save the limb.

\textbf{Contact your local emergency services for:}
\begin{itemize}
    \item A suddenly pale, cold, or painful limb, which may indicate a blocked artery.
    \item Blackened or dead-looking tissue.
    \item Spreading redness or a high fever with a wound.
    \item A large open wound or a limb that cannot be moved after injury.
    \item Rapidly worsening infection or gas under the skin, which can signal gas gangrene.
\end{itemize}

\section*{Diagnosis and Investigations}
Investigations aim to confirm the diagnosis, plan the optimal level of amputation, and assess fitness for surgery:

\begin{itemize}
    \item \textbf{Vascular assessment:} examination of pulses and the ankle-brachial index, with ultrasound or CT angiography to map the blood supply.
    \item \textbf{Imaging:} X-ray, CT, or MRI to assess bone involvement, infection, and tumour extent.
    \item \textbf{Blood tests:} full blood count, inflammatory markers, kidney function, and blood sugar.
    \item \textbf{Microbiology:} wound swabs and bone samples to guide antibiotic therapy.
    \item \textbf{Angiography:} sometimes performed with a view to attempted revascularisation before committing to amputation.
    \item \textbf{Level selection:} surgical assessment to choose the lowest amputation level likely to heal, balancing healing and function.
    \item \textbf{Pre-operative assessment:} evaluation of fitness for anaesthesia and surgery.
\end{itemize}

\section*{Management}
Management spans surgery, wound care, prosthetics, and rehabilitation:

\begin{itemize}
    \item \textbf{Timing:} emergency amputation for life-threatening infection or bleeding, and planned amputation when revascularisation has failed or is not possible.
    \item \textbf{Level:} toe, foot, below-knee, or above-knee amputation, selected to maximise healing and functional outcome.
    \item \textbf{Perioperative care:} antibiotics, venous thromboembolism prophylaxis, pain relief, and early mobilisation.
    \item \textbf{Stump care:} wound management, shaping of the stump, and prevention of joint contractures.
    \item \textbf{Prosthetics:} a custom-made artificial limb fitted once the stump has healed and settled.
    \item \textbf{Rehabilitation:} physiotherapy, occupational therapy, and gait retraining, ideally through a specialist prosthetic rehabilitation service.
    \item \textbf{Pain management:} treatment of surgical pain and later phantom limb pain with a range of pharmacological and non-pharmacological approaches.
    \item \textbf{Psychological support:} counselling and peer support for adjustment, body image, and mood.
    \item \textbf{Managing the underlying cause:} glycaemic control, smoking cessation, and treatment of vascular risk factors.
\end{itemize}

\section*{Complications}
Complications can affect the stump, the remaining limb, and general health:

\begin{itemize}
    \item Wound infection, wound breakdown, and the need for a higher-level amputation.
    \item Bleeding and haematoma formation.
    \item Phantom limb sensations and phantom limb pain, which are common and often improve with time.
    \item Stump pain and painful neuromas.
    \item Joint stiffness and contractures.
    \item Deep vein thrombosis and pulmonary embolism.
    \item Pressure sores on the stump and on the remaining leg.
    \item Depression and adjustment difficulties.
    \item Reduced mobility and falls, especially in older people.
\end{itemize}

\section*{Prognosis and Outcomes}
Most people recover well from amputation, and the majority of lower-limb amputees become mobile again, many with a prosthetic limb. Outcomes depend on the level of amputation, the person's general health, and the quality of rehabilitation; below-knee amputations generally have better walking outcomes than above-knee amputations. Phantom limb pain affects many people but usually eases over one to two years. People who lose a leg from vascular disease are at higher risk of heart attack, stroke, and amputation of the other leg, so tight control of risk factors is essential for long-term health.

\section*{Prevention and Self-Care}
Most amputations are caused by conditions that can be prevented or managed:

\begin{itemize}
    \item Stop smoking, the biggest modifiable risk factor for vascular amputation.
    \item Control diabetes tightly, including daily foot checks for early wounds.
    \item Manage blood pressure and cholesterol.
    \item Wear properly fitted footwear and see a podiatrist regularly if you have diabetes.
    \item Seek early review of any non-healing foot ulcer.
    \item Protect the stump: keep it clean, inspect it daily, and follow the rehabilitation programme.
    \item Attend physiotherapy and prosthetic appointments.
    \item Use walking aids and modify the home to prevent falls.
    \item Seek support for emotional adjustment and phantom limb pain.
\end{itemize}

\section*{Sources and Further Reading}
The following reputable sources provide reliable, up-to-date information on amputation:

\begin{itemize}
    \item NHS. \textit{Amputation information for patients.} https://www.nhs.uk/
    \item Mayo Clinic. \textit{Amputation and rehabilitation information.} https://www.mayoclinic.org/
    \item MSD Manual. \textit{Amputation overview.} https://www.msdmanuals.com/
    \item MedlinePlus. \textit{Amputation health information.} https://medlineplus.gov/
    \item Centers for Disease Control and Prevention. \textit{Diabetes complications and amputation prevention.} https://www.cdc.gov/
    \item World Health Organization. \textit{Information on disability and rehabilitation.} https://www.who.int/
\end{itemize}
